% Generated draft: requires mathematical and source review.
\documentclass[11pt]{article}
\usepackage[margin=25mm]{geometry}
\usepackage{fontspec}
\setmainfont{Latin Modern Roman}
\usepackage{amsmath,amssymb,mathtools,unicode-math}
\setmathfont{Latin Modern Math}
\usepackage{enumitem,needspace,xurl,hyperref}
\setlength{\parindent}{0pt}
\setlength{\parskip}{5pt}
\setlength{\emergencystretch}{3em}
\newcommand{\N}{\mathbb{N}}
\newcommand{\Z}{\mathbb{Z}}
\newcommand{\Q}{\mathbb{Q}}
\newcommand{\R}{\mathbb{R}}
\newcommand{\C}{\mathbb{C}}
\newcommand{\F}{\mathbb{F}}
\newcommand{\A}{\mathbb{A}}
\newcommand{\PP}{\mathbb{P}}
\newcommand{\E}{\mathbb{E}}
\newcommand{\eps}{\varepsilon}
\newcommand{\dd}{\,\mathrm{d}}
\begin{document}
\section*{UnsolvedMath ID: 20001660}
\textbf{Source problem code:} catalogue entry
\subsection{Definitions and mathematical statement}
Let \(\Gamma\) be finitely generated with word length \(|\cdot|\) and ball \(B(r)\). Property RD means there is a polynomial bound \(P\) such that every finitely supported \(f:\Gamma\to\C\) supported in \(B(r)\) satisfies \(\|\lambda(f)\|\le P(r)\|f\|_2\), where \(\lambda(f)\) is left convolution on \(\ell^2(\Gamma)\). Equivalently the multiplication form
\[
 T(f,g,h)=\sum_{abc=1}f(a)g(b)h(c)
\]
has \(|T(f,g,h)|\le P(r)\|f\|_2\|g\|_2\|h\|_2\) under that support restriction on \(f\).

This extracted record contains several workshop questions, all retained: (1) is RD invariant under quasi-isometry of finitely generated groups, including the special case of isometric Cayley graphs? (2) what geometric consequences follow from RD, for example for spaces on which a group acts cocompactly? (3) define a suitable RD notion for bounded-degree graphs or uniformly locally finite metric spaces, meaning \(\sup_x|B(x,r)|<\infty\) for every fixed \(r\); (4) Lafforgue's question asks whether RD is equivalent to a polynomial finite-set estimate
\[
 |\{(a,b,c)\in A\times B\times C:abc=1\}|
 \le P(r)\sqrt{|A||B||C|},\qquad A\subseteq B(r),
\]
for all finite subsets \(A,B,C\subseteq\Gamma\) and all radii; (5) is there a geometric characterization through a boundary or compact-space action? Item (6) is blank in the supplied extraction. Questions (2), (3), and (5) are programs with unspecified target conclusions or definitions, not a single theorem silently chosen here. RD implies (4) by indicators; the aid does not prove its converse.

For its partial result define a Lorentz norm with the explicit normalization
\[
 \|u\|_{2,1}=\int_0^\infty
       \sqrt{|\{x:|u(x)|>t\}|}\,dt
 =\sum_{j=1}^{m}(\sqrt j-\sqrt{j-1})u^*(j)
\]
for finite support of size \(m\), where \(u^*\) is the decreasing rearrangement of absolute values. Thus \(\|1_A\|_{2,1}=\sqrt{|A|}\). The aid reports that the finite-set estimate is equivalent, with the same polynomial, to
\[
 |T(f,g,h)|\le P(r)\|f\|_{2,1}\|g\|_{2,1}\|h\|_{2,1},
 \qquad\operatorname{supp}(f)\subseteq B(r).
\]
It further reports the finite-support bound
\[
 \|u\|_{2,1}\le\left(1+\frac{H_{m-1}}4\right)^{1/2}\|u\|_2,
 \qquad H_j=\sum_{i=1}^j\frac1i,\ H_0=0.
\]
The zero function is handled separately. Since finite-generation bounds ball size exponentially, the short factor incurs only \(O(\sqrt{r+1})\), or \(O(\sqrt r)\) for \(r\ge1\). The two unrestricted-support factors still have unbounded logarithmic losses; thus this does not give RD.

The aid reports a positive rank-one trilinear family with uniformly bounded indicator-set constants but strong \(\ell^2\) operator norm \(H_n^{3/2}\). Its coefficients and normalization are not supplied, so it is retained as an existence claim without an invented formula. It shows, according to the aid, that positivity, layer cake, dyadic decomposition, polarization, or random signs alone cannot upgrade a generic indicator estimate to a uniform strong \(\ell^2\) estimate. It is not a counterexample from group multiplication; a proof of the group-specific converse must use additional structure.

The original background credits de la Harpe [Ha] with RD for hyperbolic groups and Drutu--Sapir [DS2] with an extension to ``(**)-relatively hyperbolic'' groups relative to RD subgroups; the double-star condition and references are not defined here. A dated remark says the only then-known obstruction for finitely generated groups was an amenable subgroup of non-polynomial growth, and lists D. Wise's non-Hopfian groups and cocompact lattices in semisimple Lie groups as possible sources of new counterexamples. These are workshop assessments and candidates, not a current classification or verified counterexamples. The appended ``7 Applications of RD'' is an extraction heading, not a solved mathematical item.
\subsection{Short English statement}
This record bundles several RD questions. The aid addresses Lafforgue's finite-set criterion: it is equivalent to a normalized Lorentz trilinear bound, but the support-dependent losses do not establish strong \(\ell^2\) RD. A generic positive-form example explains why additional group structure is needed.
\subsection{Sources}
\begin{itemize}
\item Local UnsolvedMath registry, permanent ID \(20001660\), source code \texttt{catalogue entry}; \texttt{aim-geometric-group-theory-notes.json}, \texttt{notes[168]}, tag \texttt{problem}.
\item AIM workshop document ``Property of rapid decay,'' dated 21 March 2006, source item 1: \url{https://aimath.org/WWN/rapiddecay/rapiddecay.pdf}. Source type PDF; statement status \texttt{exact}. Supplied verification reports an inspection of official PDF page 5 and notes joined headings/remarks. That inspection is source metadata; no PDF was accessed or compiled for this draft.
\item Original labels [Ha], P. de la Harpe, and [DS2], C. Drutu and M. Sapir; Lafforgue's finite-set question; D. Wise's candidate groups. No full bibliographic entries, authoritative citations for the Lorentz endpoint audit, or explicit rank-one family coefficients are supplied. Original literature field is empty.
\item Research attempt 1: \texttt{partial\_result}, status at run \texttt{partially\_solved}, rigor label \texttt{proved}, novelty confidence low; registry status \texttt{partially\_solved}, classification \texttt{PARTIAL-PROGRESS}. Created and updated 14 August 2026. Machine-generated aid, not peer reviewed; independently verify before citation.
\item Ulam AI / UnsolvedMath Contributors; CC BY 4.0.
\end{itemize}
\end{document}
